\section{Actual cell displacement and time reversal}\label{sec:current}
For $z\in M$ outside the singularity set, lift its current collision to the
disk at the origin.  If the next collision is on the disk at
$\kappa_R(z)=m_1e_1+m_2e_2$, define
\[
 \kappa_R(z)=(m_1,m_2)\in\Z^2,
 \qquad J_R(z)=m_1.
\]
Finite horizon makes $\kappa_R$ and $J_R$ bounded and piecewise H\"older.
Let $\mathscr R(\alpha,\varphi)=(\alpha,-\varphi)$ denote velocity reversal.
Then
\[
 T_R^{-1}=\mathscr R T_R\mathscr R,
 \qquad J_R\circ\mathscr R\circ T_R=-J_R.
\]
Consequently $\int J_R\,d\mu=0$.

\begin{proposition}[Open displacement branches and lattice span]\label{prop:span}
For every $R\in\mathcal I$, there are open continuity components of the
collision map on which
\[
 \kappa_R=(1,0)
 \qquad\text{and}\qquad
 \kappa_R=(0,1),
\]
respectively.  Their time reversals have the opposite displacements.  Hence the
support of the actual displacement cocycle generates $\Z^2$, and the scalar
current $J_R$ has arithmetic span one.
\end{proposition}
\begin{proof}
The normal free flight from the disk at the origin to its nearest neighbour at
$e_1$ has length $1-2R>0$.  Its centre line stays at distance at least
$\sqrt3/2>R$ from every other lattice centre.  Therefore a small open set of
nearby outgoing collision states also hits the $e_1$ disk first, giving
$\kappa_R=(1,0)$.  The identical argument in the $e_2$ direction gives
$\kappa_R=(0,1)$.  Velocity reversal gives the negative branches.  The four
values generate the full lattice, and the first coordinate takes adjacent
integer values zero and one.
\end{proof}

\section{Nonconjugate radius response}\label{sec:nonconjugacy}
There is a normal two-bounce orbit between neighbouring disks at separation
one, with period
\[
 L_1(R)=2(1-2R).
\]
There is also a normal two-bounce orbit between copies separated by
$e_1+e_2$, whose centre distance is $\sqrt3$.  The line passes at distance
$1/2$ from the intermediate lattice centre, and
$R\le47/100<1/2$, so the orbit is unobstructed.  Its period is
\[
 L_{\sqrt3}(R)=2(\sqrt3-2R).
\]

\begin{theorem}[No global similarity or constant time rescaling]\label{thm:nonconjugacy}
If $R\ne R_0$, the billiard flows at radii $R$ and $R_0$ are not related by a
lattice-preserving conjugacy with one constant time factor.
\end{theorem}
\begin{proof}
Such a conjugacy would preserve the ratio of all marked periods.  But
\[
 \frac{L_{\sqrt3}(R)}{L_1(R)}
 =\frac{\sqrt3-2R}{1-2R}
\]
is strictly increasing in $R$.
\end{proof}

\begin{proposition}[Explicit mean-roof response]\label{prop:santalo}
The Liouville mean free-flight time is
\[
 \bar\tau_R^{\rm L}
 =\int_M\tau_R\,d\mu
 =\frac{A_\Lambda-\pi R^2}{2R},
\]
and
\[
 \frac{d}{dR}\bar\tau_R^{\rm L}
 =-\frac{A_\Lambda}{2R^2}-\frac\pi2<0.
\]
\end{proposition}
\begin{proof}
Santal\'o's return-time formula gives
$\pi\abs{Q_R}/\abs{\partial Q_R}$.  Substitute
$\abs{Q_R}=A_\Lambda-\pi R^2$ and
$\abs{\partial Q_R}=2\pi R$.
\end{proof}

\section{The twisted Lorentz spectral packet}\label{sec:spectral}
Let $\cL_R$ be the Perron--Frobenius operator of $T_R$ relative to $\mu$ and
set
\[
 \cL_{R,q}f=\cL_R(e^{qJ_R}f).
\]
The following packet is the precise external input from finite-horizon Lorentz
spectral theory.

\begin{theorem}[Uniform twisted spectral packet]\label{thm:spectral-packet}
There are $q_*>0$, $0<\rho_*<1$, and anisotropic Banach spaces
$\mathcal B\hookrightarrow\mathcal B_w$, valid uniformly for
$R\in\mathcal I$, such that for $\abs q<q_*$,
\[
 \cL_{R,q}^n=e^{n\Lambda_R(q)}\Pi_{R,q}+N_{R,q}^n,
 \qquad \norm{N_{R,q}^n}_{\mathcal B\to\mathcal B_w}\le C\rho_*^n.
\]
The leading eigenvalue is simple, positive for real $q$, and analytic in $q$.
The spectral data vary continuously with $R$.  The same packet holds for
complex $q+it$ near the real axis and gives the lattice local limit theorem
for $J_R$.
\end{theorem}
\begin{proof}
Finite horizon, strictly positive curvature, bounded complexity, and the
uniform radius interval give the common hyperbolicity and homogeneity-strip
constants required by the Lorentz anisotropic spaces.  Quasicompactness,
large deviations, and ASIP follow from \cite{DemersZhang2011}; persistence
